% !TEX spellcheck = en_US
% !TEX spellcheck = LaTeX
\documentclass[letterpaper,10pt,english]{article}
\usepackage{%
amsfonts,%
amsmath,%
amssymb,%
amsthm,%
babel,%
bbm,%
%biblatex,%
caption,%
centernot,%
color,%
enumerate,%
%enumitem,%
epsfig,%
epstopdf,%
etex,%
fancybox,%
framed,%
fullpage,%
%geometry,%
graphicx,%
hyperref,%
latexsym,%
mathptmx,%
mathtools,%
multicol,%
paralist,%
pgf,%
pgfplots,%
pgfplotstable,%
pgfpages,%
proof,%
psfrag,%
%subfigure,%
tcolorbox,%
tfrupee,%
tikz,%
times,%
%ulem,%
url,%
xcolor,%
mathpazo,%
}

\definecolor{shadecolor}{gray}{.95}%{rgb}{1,0,0}
\usepackage[margin=1in,top=0.75in]{geometry}
\usepackage[mathscr]{eucal}
\usepgflibrary{shapes}
\usepgfplotslibrary{fillbetween}
\usetikzlibrary{%
arrows,%
backgrounds,%
chains,%
decorations.pathmorphing,% /pgf/decoration/random steps | erste Graphik
decorations.text,% 
matrix,%
positioning,% wg. " of "
fit,%
patterns,%
petri,%
plotmarks,%
scopes,%
shadows,%
shapes.misc,% wg. rounded rectangle
shapes.arrows,%
shapes.callouts,%
shapes%
}

%\pgfplotsset{compat=newest} %<------ Here
\pgfplotsset{compat=1.11} %<------ Or use this one

\theoremstyle{plain}
\newtheorem{thm}{Theorem}[section]
\newtheorem{lem}[thm]{Lemma}
\newtheorem{prop}[thm]{Proposition}
\newtheorem{cor}[thm]{Corollary}
\newtheorem{clm}[thm]{Claim}

\theoremstyle{definition}
\newtheorem{defn}[thm]{Definition}
\newtheorem{conj}[thm]{Conjecture}
\newtheorem{exmp}[thm]{Example}
\newtheorem{axiom}[thm]{Axiom}
\newtheorem{assum}[thm]{Assumption}
\newtheorem{exerc}[thm]{Exercise}
\newtheorem*{defin}{Definition}

\theoremstyle{remark}
\newtheorem{rem}{Remark}
\newtheorem{note}{Note}

\newcommand{\iid}{\emph{i.i.d.}\ }
\newcommand{\Cov}{\operatorname{Cov}}
%\newcommand{\det}{\operatorname{det}}
%\newcommand{\Real}{\mathbb{R}}
\newcommand{\tr}{\operatorname{tr}}
\newcommand{\Var}{\operatorname{Var}}
\newcommand{\cov}{\operatorname{cov}}

%\renewcommand{\proof}[1]{\begin{proof}#1\end{proof}}
\newcommand{\EQ}[1]{\begin{equation*}#1\end{equation*}}
\newcommand{\EQN}[1]{\begin{equation}#1\end{equation}}
\newcommand{\eq}[1]{\begin{align*}#1\end{align*}}
\newcommand{\meq}[2]{\begin{xalignat*}{#1}#2\end{xalignat*}}
\newcommand{\norm}[1]{\left\lVert#1\right\rVert}
\newcommand{\inner}[1]{\left\langle #1\right\rangle}
\newcommand{\abs}[1]{\left\lvert#1\right\rvert}
\newcommand{\expect}[1]{\mathbb{E}\left[{#1}\right]}
\newcommand{\prob}[1]{\mathbb{P}\left[{#1}\right]}
\newcommand{\given}{\; \big\vert \;} 
\newcommand{\set}[1]{\left\{#1\right\}} 
\newcommand{\indicator}[1]{1_{\set{#1}}} 
\newcommand{\Ind}[1]{\mathbbm{1}_{#1}} 
\newcommand{\SetIn}[1]{\mathbbm{1}_{\set{#1}}} 
\newcommand{\red}[1]{\textcolor{red}{#1}} 
\newcommand{\floor}[1]{\left\lfloor#1\right\rfloor}
\newcommand{\ceil}[1]{\left\lceil#1\right\rceil}


\newcommand{\C}{\mathbb{C}}
\newcommand{\D}{\mathbb{D}}
\newcommand{\E}{\mathbb{E}}
\newcommand{\F}{\mathbb{F}}
\newcommand{\N}{\mathbb{N}}
\renewcommand{\P}{\mathbb{P}}
\newcommand{\Q}{\mathbb{Q}}
\newcommand{\R}{\mathbb{R}}
\newcommand{\Z}{\mathbb{Z}}

\newcommand{\bA}{\mathbf{A}}
\newcommand{\bI}{\mathbf{I}}
\newcommand{\bU}{\mathbf{1}}
\newcommand{\bx}{\mathbf{x}}
\newcommand{\bX}{\mathbf{X}}
\newcommand{\bZ}{\mathbf{Z}}


\newcommand{\cA}{\mathcal{A}}
\newcommand{\cB}{\mathcal{B}}
\newcommand{\cC}{\mathcal{C}}
\newcommand{\cF}{\mathcal{F}}
\newcommand{\cG}{\mathcal{G}}
\newcommand{\cM}{\mathcal{M}}
\newcommand{\cN}{\mathcal{N}}
\newcommand{\cP}{\mathcal{P}}
\newcommand{\cT}{\mathcal{T}}
\newcommand{\cX}{\mathcal{X}}
\newcommand{\cY}{\mathcal{Y}}

\newcommand{\sA}{\mathscr{A}}
\newcommand{\sB}{\mathscr{B}}
\newcommand{\sC}{\mathscr{C}}
\newcommand{\sE}{\mathscr{E}}
\newcommand{\sF}{\mathscr{F}}
\newcommand{\sG}{\mathscr{G}}
\newcommand{\sH}{\mathscr{H}}
\newcommand{\sL}{\mathscr{L}}
\newcommand{\sS}{\mathscr{S}}
\newcommand{\sT}{\mathscr{T}}
\newcommand{\sX}{\mathscr{X}}
\newcommand{\sY}{\mathscr{Y}}
\newcommand{\sZ}{\mathscr{Z}}

\newcommand{\tS}{\tilde{S}}
% Debug
\newcommand{\todo}[1]{\begin{color}{blue}{{\bf~[TODO:~#1]}}\end{color}}

% a few handy macros

\renewcommand{\le}{\leqslant}
\renewcommand{\ge}{\geqslant}
\newcommand\matlab{{\sc matlab}}
\newcommand{\goto}{\rightarrow}
\newcommand{\bigo}{{\mathcal O}}
%\newcommand{\half}{\frac{1}{2}}
%\newcommand\implies{\quad\Longrightarrow\quad}
\newcommand\reals{{{\rm l} \kern -.15em {\rm R} }}
\newcommand\complex{{\raisebox{.043ex}{\rule{0.07em}{1.56ex}} \hskip -.35em {\rm C}}}


% macros for matrices/vectors:

% matrix environment for vectors or matrices where elements are centered
\newenvironment{mat}{\left[\begin{array}{ccccccccccccccc}}{\end{array}\right]}
\newcommand\bcm{\begin{mat}}
\newcommand\ecm{\end{mat}}

% matrix environment for vectors or matrices where elements are right justifvied
\newenvironment{rmat}{\left[\begin{array}{rrrrrrrrrrrrr}}{\end{array}\right]}
\newcommand\brm{\begin{rmat}}
\newcommand\erm{\end{rmat}}

% for left brace and a set of choices
%\newenvironment{choices}{\left\{ \begin{array}{ll}}{\end{array}\right.}
\newcommand\when{&\text{if~}}
\newcommand\otherwise{&\text{otherwise}}
% sample usage:
%\delta_{ij} = \begin{choices} 1 \when i=j, \\ 0 \otherwise \end{choices}


% for labeling and referencing equations:
\newcommand{\eql}{\begin{equation}\label}
\newcommand{\eqn}[1]{(\ref{#1})}
% can then do
%\eql{eqnlabel}
%...
%\end{equation}
% and refer to it as equation \eqn{eqnlabel}.
% some useful macros for finite difference methods:
\newcommand\unp{U^{n+1}}
\newcommand\unm{U^{n-1}}
% for chemical reactions:
\newcommand{\react}[1]{\stackrel{K_{#1}}{\rightarrow}}
\newcommand{\reactb}[2]{\stackrel{K_{#1}}{~\stackrel{\rightleftharpoons}
 {\scriptstyle K_{#2}}}~}
\makeatletter
\def\th@plain{%
\thm@notefont{}% same as heading font
\itshape % body font
}
\def\th@definition{%
\thm@notefont{}% same as heading font
\normalfont % body font
}
\makeatother
\date{}


\title{Lecture-08: Expectation}
\author{}

\begin{document}
\maketitle

\section{Expectation}
%Consider a set of outcomes $\Omega$, the set of events $\sF$, and the probability function $P: \sF \to [0,1]$. 
\begin{shaded}
\begin{exmp}
Consider a probability space $(\Omega, \sF, P)$.
We consider $N$ trials of a random experiment, 
and define a random vector $X: \Omega \to \sX^N$ 
such that $X_i \triangleq \pi_i\circ X :\Omega \to\sX$ is a discrete random variable associated with the trial $i \in [N]$.   
If the marginal distributions of random variables $(X_1, \dots, X_N)$ are identical with the common probability mass function $P_{X_1}: \sX \to [0,1]$,  
then the empirical mean of random variable $X_1$ can be written as 
\EQ{
\hat{m} \triangleq \frac{1}{N}\sum_{i=1}^NX_i. 
}
%If $X: \Omega \to \sX \subseteq \R$ is a discrete random variable with probability mass function (PMF) $P: \sX \to [0,1]$,  
For a random variable $X_1: \Omega \to \sX$, 
we can define events $E_{X_1}(x) \triangleq X_1^{-1}\set{x}$ for each value $x \in \sX$.  
The probability mass function $P_{X_1}:\sX \to [0,1]$ for the discrete random variable $X_1$ can be estimated for each $x \in \sX$ as the empirical probability mass function
\EQ{
\hat{P}_{X_1}(x) \triangleq \frac{1}{N}\sum_{i=1}^N\Ind{E_{X_i}(x)}. 
}

Recall that a simple random variable $X_1$ can be written as $X_1 = \sum_{x \in \sX}x\Ind{E_{X_1}(x)}$, 
where $E_{X_1} \triangleq (E_{X_1}(x) \in \sF: x \in \sX)$ is a finite partition of the sample space $\Omega$ and $P_{X_1}(x) = P(E_{X_1}(x))$. 
That is, we can write the empirical mean in terms of the empirical PMF as 
\EQ{
\hat{m} 
= \frac{1}{N}\sum_{i=1}^N\sum_{x \in \sX}x\Ind{E_{X_i}(x)}(\omega)
= \sum_{x \in \sX}x\hat{P}_{X_1}(x).
}
\end{exmp}
\end{shaded}

This example motivates the following definition of mean for simple random variables. 
\begin{defn}[Expectation of simple random variable] 
%A discrete random variable $X: \Omega \to \sX \subseteq \R$ defined on a probability space $(\Omega, \sF, P)$,  taking finitely many values $\sX$, and having probability mass function $P_X: \sX \to [0,1]$ is called a \textbf{simple random variable}.  
The \textbf{mean} or \textbf{expectation} of a simple random variable $X:\Omega \to \sX \subseteq \R$ defined on a probability space $(\Omega, \sF, P)$, is denoted by $\E[X]$ and defined as 
\EQ{
\E[X] \triangleq \sum_{x \in \sX}xP_X(x).
}
\end{defn}
\begin{rem} 
For an indicator random variable $\Ind{A}$, we have $\E\Ind{A} = P(A)$. 
That is, the expectation of an indicator function is the probability of the indicated set. 
\end{rem}
\begin{rem} 
%For a random variable $X: \Omega \to \sX$, we can define events $E_X(x) \triangleq X^{-1}\set{x}$ for each value $x \in \sX$.  
%Recall that a simple random variable can be written as $X = \sum_{x \in \sX}x\SetIn{X = x}$, where $E_X \triangleq (E_X(x) \in \sF: x \in \sX)$ is a finite partition of the sample space $\Omega$ and $P_X(x) = P(E_X(x))$. 
Since a simple random variable can be written as $X = \sum_{x \in \sX}x\Ind{E_X(x)}$ where  $E_X(x) \triangleq X^{-1}\set{x}$ for all $x\in\sX$, 
we can write the expectation of a simple random variable as an integral 
\EQ{
\E[X] 
= \int_{\Omega}X(\omega)P(d\omega) 
= \int_\Omega\sum_{x \in \sX}x\Ind{E_X(x)}(\omega)P(d\omega) 
= \sum_{x \in \sX}x\int_{\Omega}\Ind{E_X(x)}(\omega)P(d\omega) 
= \sum_{x \in \sX}x\E[\Ind{E_X(x)}] = \sum_{x \in \sX}xP_X(x).
}
\end{rem}
\begin{thm}
Consider a non-negative random variable $X: \Omega \to \R_+$ defined on a probability space $(\Omega, \sF, P)$.  
There exists a sequence of non-decreasing non-negative simple random variables $Y: \Omega \to \R_+^\N$ such that for all $\omega \in \Omega$
\EQ{
Y_n(\omega) \le Y_{n+1}(\omega), \text{ for all } n \in \N, \text{ and  } \lim_{n}Y_n(\omega) = X(\omega).
}
Then $\E[Y_n]$ is defined for each $n \in \N$, the sequence $(\E[Y_n] \in \R_+: n \in \N)$ is non-decreasing, 
and the limit $\lim_n\E[Y_n] \in \R_+\cup\set{\infty}$ exists.  
This limit is independent of the choice of the sequence and depends only on the probability space. % $(\Omega, \sF, P)$.   
\end{thm}
\begin{proof}
For each $n \in \N$ and $k \in \set{0, \dots, 2^{2n}-1}$, 
we define half-open sets $B_{n,k} \triangleq (k2^{-n}, (k+1)2^{-n}]$. 
Then, the collection of sets $B_n \triangleq (B_{n,k}: k \in \set{0, \dots, 2^{2n}-1})$ partitions the set $(0, 2^n]$ for each $n \in \N$. 
Further, we observe that $\cup_{n \in \N}(0, 2^n] = \R^+$ and that $B_{n+1,2k}\cup B_{n+1,2k+1} = B_{n,k}$ for all $n \in \N$ and $k$. 

For a non-negative random variable $X: \Omega \to \R_+$, 
we define events $A^X_{n,k} = X^{-1}(B_{n,k}) \in \sF$, 
and a sequence of simple non-negative random variables $Y: \Omega\to\R_+^\N$ in the following fashion
\EQ{
Y_n(\omega) 
\triangleq \sum_{k=0}^{2^{2n}-1}\Ind{A^X_{n,k}}(\omega)\left(\inf_{\omega \in A^X_{n,k}}X(\omega)\right)
= \sum_{k=0}^{2^{2n}-1}\Ind{A^X_{n,k}}(\omega)\left(\inf_{X(\omega) \in B_{n,k}}X(\omega)\right)
= \sum_{k=0}^{2^{2n}-1}k2^{-n}\Ind{A^X_{n,k}}(\omega).
%= \sum_{k=0}^{2^{2n}-1}k2^{-n}\SetIn{[k2^{-n}, (k+1)2^{-n})}(\omega)
%= \begin{cases}
%k2^{-n}, & k2^{-n} < X(\omega) \le (k+1)2^{-n}, k \in \set{0, \dots, 2^{2n}-1},\\
%0, & X(\omega) > 2^n.
%\end{cases}
}
We observe that $Y_n$ is a quantized version of $X$, 
and its value is the left end-point $k2^{-n}$ when $X \in B_{n,k}$ for each $k \in \set{0, \dots, 2^{2n}-1}$.  
Since $\cup_{k=0}^{2^{2n}-1}A^X_{n,k} = X^{-1}(0, 2^n]$, 
it follows that we cover the positive real line as $n$ grows larger and the step size grows smaller. 
Thus, the limiting random variable can take all possible non-negative real values. 
We observe that 
\eq{
Y_{n+1}(\omega) 
&= \sum_{k=0}^{2^{2(n+1)}-1}\Ind{A^X_{n+1,k}}(\omega)\left(\inf_{X(\omega) \in B_{n+1,k}}X(\omega)\right)\\
&\ge \sum_{k=0}^{2^{2n}-1}\left(\Ind{A^X_{n+1,2k}}(\omega)\left(\inf_{X(\omega) \in B_{n+1,2k}}X(\omega)\right)
+ \Ind{A^X_{n+1,2k+1}}(\omega)\left(\inf_{X(\omega) \in B_{n+1,2k+1}}X(\omega)\right)\right)\\
&\ge  \sum_{k=0}^{2^{2n}-1}\Ind{A^X_{n,2k}}(\omega)\left(\inf_{X(\omega) \in B_{n,k}}X(\omega)\right)
= Y_n(\omega).
}
We see that $Y_n(\omega) \le Y_{n+1}(\omega) \le X(\omega)$ and $\lim_nY_n(\omega) = X(\omega)$ for all $\omega \in \Omega$. 

Since $Y_n: \Omega \to \R$ is a simple random variable for all $n \in \N$, 
the expectation $\E[Y_n]$ is defined for all $n$, 
and can be written as 
\EQ{
\E[Y_n]
= \sum_{k=0}^{2^{2n}-1}k2^{-n}[F_X((k+1)2^{-n})- F_X(k2^{-n})].
}
We observe that this expectation is completely specified by the distribution function $F_X$, 
and we can write the limit 
\EQ{
\lim_n\E[Y_n] 
= \lim_n\sum_{k=0}^{2^{2n}-1}k2^{-n}[F_X(k2^{-n}+2^{-n})- F_X(k2^{-n)})] 
= \int_{\R^+} x dF_X(x).
}
\end{proof}
\begin{defn}[Expectation of a non-negative random variable] 
For a non-negative random variable $X: \Omega \to \R$ defined on the probability space $(\Omega,\sF,P)$, 
consider the sequence of non-decreasing simple random variables $Y: \Omega \to \R_+^\N$ such that $\lim_nY_n = X$. 
The \textbf{expectation} of the non-negative random variable $X$ is defined as 
\EQ{
\E[X] \triangleq \lim_n\E[Y_n]. 
}
\end{defn}
\begin{rem} 
From the definition, it follows that $\E[X] = \int_{\R_+} x dF_X(x)$. 
\end{rem}
\begin{defn}[Expectation of a real random variable] 
For a real-valued random variable $X$ defined on a probability space $(\Omega, \sF, P)$, 
we can define the following functions
\begin{xalignat*}{2}
&X_+ \triangleq \max\set{X, 0}, &
&X_- \triangleq \max\set{0, -X}.
\end{xalignat*}
We can verify that $X_+, X_-$ are non-negative random variables and hence their expectations are well defined.
We observe that $X(\omega) = X_+(\omega) - X_-(\omega)$ for each $\omega \in \Omega$. 
If at least one of the $\E[X_+]$ and $\E[X_-]$ is finite, then the \textbf{expectation} of the random variable $X$ is defined as 
\EQ{
\E[X] \triangleq \E[X_+] - \E[X_-]. 
}
\end{defn}
\begin{thm}[Expectation as an integral with respect to the distribution function] 
For a random variable $X:\Omega\to\R$ defined on the probability space $(\Omega, \sF, P)$, 
the expectation is given by 
\EQ{
\E[X] = \int_\R x dF_X(x).
}
\end{thm}
\begin{proof}
It suffices to show this for a non-negative random variable $X$, 
and the result follows from the definition of expectation of a non-negative random variable 
as the limit of expectation of approximating simple functions.
%Recall that there exists a sequence of non-decreasing simple random variables $Y: \Omega \to \R_+^\N$, 
%such that we can write the expectation for each term as 
%\EQ{
%\E[Y_n] 
%= \sum_{k=0}^{2^{2n}-1}k2^{-n}[F_X((k+1)2^{-n})- F_X(k2^{-n})],
%}
%is completely specified by the distribution function $F_X$. 
%It follows that the expectation of the random variable $X$ is given by
%\EQ{
%\E[X] 
%= \lim_n\E[Y_n] 
%= \lim_n\sum_{k=0}^{2^{2n}-1}k2^{-n}[F_X(k2^{-n}+2^{-n})- F_X(k2^{-n)})] 
%= \int_{\R^+} x dF_X(x).
%}
\end{proof}                                                                                                     
\section{Properties of Expectations}
\begin{thm}[Properties] 
Let $X: \Omega \to \R$ be a random variable defined on the probability space $(\Omega, \sF, P)$. 
\begin{enumerate}[(i)]
\item \textbf{Linearity:} Let $a, b \in \R$ and $X, Y$ be random variables defined on the probability space $(\Omega, \sF, P)$. 
If $\E X, \E Y$, and $a\E X + b\E Y$ are well defined, then $\E(aX+bY)$ is well defined and 
\EQ{
\E(aX+bY) = a\E X + b \E Y.
}
\item \textbf{Monotonicity:} If $P\set{X \ge Y} = 1$ and $\E[Y]$ is well defined
with $\E[Y] > -\infty$, then $\E[X]$ is well defined and $\E[X] \ge \E[Y]$. 
\item \textbf{Functions of random variables:} Let $g: \R \to \R$ be a Borel measurable function, then $g(X)$ is a random variable with $\E[g(X)] = \int_{x\in \R}g(x)dF(x)$. 
\item \textbf{Continuous random variables:} Let $f_X: \R \to [0, \infty)$ be the density function, 
then $\E X = \int_{x \in \R}xf_X(x)dx$. 
\item \textbf{Discrete random variables:} Let $P_X: \sX \to [0,1]$ be the probability mass function, 
then $\E X = \sum_{x \in \sX}xP_X(x)$. 
\item \textbf{Integration by parts:} The expectation $\E X = \int_{x \ge 0}(1-F_X(x))dx - \int_{x < 0}F_X(x)dx$ is well defined when at least one of the two parts is finite on the right hand side. 
\end{enumerate}
\end{thm}
\begin{proof}
It suffices to show properties $(i)-(iii)$ for simple random variables. 
\begin{enumerate}[(i)]
\item Let $X = \sum_{x \in \sX}x\mathbbm{1}_{E_X(x)}$ and $Y =  \sum_{y \in \sY}y\mathbbm{1}_{E_Y(y)}$ be simple random variables, then $(E_X(x)\cap E_Y(y) \in \sF: (x, y)\in \sX \times \sY)$ partition the sample space $\Omega$. 
Hence, we can write $aX + bY = \sum_{(x,y) \in \sX \times \sY}(ax+by)\SetIn{E_X(x)\cap E_Y(y)}$ and from linearity of sum it follows that
\eq{
\E[aX+bY] &= \sum_{(x,y) \in \sX \times \sY}(ax+by)P\set{E_X(x)\cap E_Y(y)} \\
&= a\sum_{x \in \sX}x\sum_{y\in\sY}P\set{E_X(x)\cap E_Y(y)} 
+ b\sum_{y \in \sY}y\sum_{x \in \sX}P\set{E_X(x)\cap E_Y(y)}\\
&= a\sum_{x \in \sX}xP(E_X(x)) + b\sum_{y \in \sY}yP(E_Y(y)) = a\E X + b \E Y.
}
\item From the fact that $X-Y \ge 0$ almost surely and linearity of expectation, 
it suffices to show that $\E X \ge 0$ for non-negative random variable $X$. 
It can easily be shown for simple non-negative random variables, 
and follows for general non-negative random variables by taking limits. 
%Suppose the event $E_+ \triangleq \set{X \ge 0} \in \sF$ has probability unity. %, and $\E X = a <0$. 
%We define $E_- \triangleq \set{X < 0}$ and $\E X \ge $

\item It suffices to show this holds true for simple random variables $X: \Omega \to \sX \subset \R$. 
Since $g: \R \to \R$ is Borel measurable, $Y \triangleq g(X): \Omega\to\sY \triangleq g(\sX)$ is a random variable. 
It follows that, we can write the following disjoint union $\sX = \cup_{y\in \sY}\cup_{x \in g^{-1}\set{y}}\set{x}$. 
Further, for each $y \in \sY$, we have 
\EQ{
E_Y(y) = \set{\omega \in \Omega: (g\circ X)(\omega) = y} = X^{-1}\circ g^{-1}\set{y} = \cup_{x\in g^{-1}\set{y}}E_X(x).
}
Since $E_X(x)$ are disjoint for all $x \in \sX$, we get $P(E_Y(y)) = \sum_{x \in g^{-1}\set{y}}P(E_X(x))$. 
Using the above two facts, we can write the expectation 
\EQ{
\E[Y] 
= \sum_{y \in \sY}yP(E_Y(y)) 
= \sum_{y \in g(\sX)}\sum_{x \in g^{-1}(y)}g(x)P(E_X(x)) = \sum_{x \in \sX}g(x)P(E_X(x)).
}
  
\item For continuous random variables, we have $dF_X(x) = f_X(x)dx$ for all $x \in \R$. 
\item For discrete random variables $X : \Omega \to \sX$, we have $dF_X(x) = P_X(x)$ for all $x \in \sX$ and zero otherwise. 
\item We can write $\E X = \int_{x <0}xdF_X(x)-\int_{x \ge 0}xd(1-F_X)(x)$. 
We apply integration by parts to the first term on the right, to get 
\EQ{
\int_{x <0}xdF_X(x) = xF_X(x)\rvert_{-\infty}^0 - \int_{x < 0}F_X(x)dx.
}
Similarly, we apply integration by parts to the second term on the right, to get 
\EQ{
-\int_{x \ge 0}xd(1-F_X)(x) = -x(1-F_X(x))\rvert_0^{\infty}+ \int_{x \ge 0}(1-F_X(x))dx.
}
\end{enumerate}
\end{proof}

\end{document}

\subsection{Linearity of expectations}
\begin{thm}[Linearity of expectations] 
Suppose $X: \Omega \to \R^n$ is a random vector defined on the probability space $(\Omega, \sF, P)$. 
Then, for constants $\alpha_i \in \R$ for all $i \in [n]$, we have
\EQ{
\E[\sum_{i=1}^n\alpha_iX_i] = \sum_{i=1}^n\alpha_i\E[X_i].
}
\end{thm}
\begin{proof}
It suffices to show that $\int_{x \in \R^n}x_1dF_{X}(x) = \int_{x_1 \in \R}x_1dF_{X_1}(x_1)$. 
The result follows from the observation that 
\EQ{
\int_{x \in \R^n}x_1dF_{X}(x) = \int_{x_1 \in \R}x_1\int_{x_2, \dots, x_n}dF_X(x) = \int_{x_1 \in \R}x_1dF_{X_1}(x_1). 
}
\end{proof}
%\subsection{Expectation of function of random variables}
\subsection{Functions of random variables}
%Consider a random variable $X: \Omega \to \R$ defined on the probability space $(\Omega, \sF, P)$. 
%Suppose $g: \R \to \R$ is function such that $g^{-1}(-\infty, x] \in \sB(\R)$, 
%then $g(X)$ is a random variable. 
%
%\begin{shaded}
%\begin{exmp}[Monotone function of random variables] 
%Let $g:\R \to \R$ be a monotonically increasing function such that $g^{-1}(-\infty, x] \in \sB(\R)$ for all $x \in \R$. 
%Consider a random variable $X: \Omega \to \R$ defined on the probability space $(\Omega, \sF, P)$, 
%then $Y \triangleq g(X)$ is a random variable with distribution function
%\EQ{
%F_Y(y) = P\set{g(X) \le y} = P\set{X \le g^{-1}(y)} = F_{X}(g^{-1}(y)). 
%}
%\end{exmp}
%\end{shaded}
%
%\begin{shaded}
%\begin{exmp}[Independence of function of random variables] 
%Let $g:\R \to \R$ and $h: \R \to \R$ be functions such that $g^{-1}(-\infty, x]$ and $h^{-1}(-\infty, x]$ are Borel sets for all $x \in \R$. 
%Consider independent random variables $X$ and $Y$ defined on the probability space $(\Omega, \sF, P)$. 
%Can you show that $g(X)$ and $h(Y)$ are independent random variables?
%\end{exmp}
%\end{shaded}
\begin{thm}[Fundamental theorem of expectations] 
Suppose $X$ is a random variable defined on the probability space $(\Omega, \sF, P)$, 
and $g: \R \to \R$ is a function such that $g^{-1}(-\infty, x] \in \sB(\R)$ for all $x \in \R$. 
Then $Y = g(X)$ is also a random variable, 
and 
\EQ{
\E[Y] = \E[g(X)].
}
\end{thm}
\begin{proof}
It suffices to show this is true for simple random variables $X: \Omega \to \sX \subseteq \R$. 
Let $E_X(x) \triangleq \set{\omega \in \Omega: X(\omega) = x}$ for each $x \in \sX$. 
Then $(E_X(x) = X^{-1}\set{x} \in \sF: x \in \sX)$ partitions the sample space $\Omega$, 
and we can write its expectation as 
\EQ{
\E[X] =\E[\sum_{x \in \sX}x\SetIn{E_X(x)}] = \sum_{x \in \sX}xP_X(x).
}
It follows that $Y: \Omega \to \sY = g(\sX)$ is also a discrete random variable, and we can write 
\EQ{
E_Y(y) \triangleq \set{\omega \in \Omega: (g\circ X)(\omega) = y} = \cup_{x \in \sX}\set{X(\omega) = x, g(x) = y} = \cup_{x \in g^{-1}\set{y}} E_X(x).
} 
Therefore, we can write the expectation
\EQ{
\E[Y] = \E[\sum_{y \in \sY}y\mathbbm{1}_{E_Y(y)}] = \E[\sum_{y \in \sY}y\sum_{x \in g^{-1}\set{y}}\mathbbm{1}_{E_X(x)}] = \E[\sum_{x \in \sX}g(x)\mathbbm{1}_{E_X(x)}]= \sum_{x \in \sX}g(x)P_X(x).
}
\end{proof}

\appendix
\section{Indicator functions}
For a sequence of disjoint events $(A_n \in \sF:n \in \N)$, we have 
\EQ{
\SetIn{\cup_{n \in \N}A_n} = \sum_{n \in \N}\mathbbm{1}_{A_n}. 
}
For any sequence of events $(A_n\in \sF: n \in \N)$, we have 
\EQ{
\SetIn{\cap_{n \in \N}A_n} = \prod_{n \in \N}\mathbbm{1}_{A_n}.
}
\subsection{Law of total probability}
Let $(A_n \in \sF: n \in \N)$ be a partition of the sample space, i.e. $A_n\cap A_m = \emptyset$ for $n \neq m$ and $\cup_{n \in \N}A_n = \Omega$. 
Then, any event $B = B\cap\Omega = B\cap(\cup_{n \in \N}A_n) = \cup_{n \in \N}(B\cap A_n)$. 
Therefore, we can write 
\EQ{
\mathbbm{1}_B = \sum_{n \in \N}\mathbbm{1}_{B\cap A_n}.
}
\end{document}
